\chapter{Monitoring \& the Run Summary}
\label{ch:monitoring}

\newthought{Observability without interference.} A running scan can be inspected from a
second terminal at any time --- \cli{Jarvis2 monitor} reads the runtime's shared state
and prints a snapshot without writing a single byte to it --- and every run closes with
a performance accounting in the console and on disk. This chapter covers both ends.

\section{\clih{Jarvis2 monitor}}
\label{sec:monitor-cli}

\begin{terminal}
\begin{jcmd}
Jarvis2 monitor
\end{jcmd}
\end{terminal}

Prints one text snapshot of the active scan and exits:

\begin{itemize}
  \item \textbf{Workers} --- alive/total, each Worker's status, current Sample, held
        calculator packs, and subprocess \textsc{pid}s;
  \item \textbf{queues} --- task queue and archive queue lengths;
  \item \textbf{samples} --- running / completed / failed counters;
  \item \textbf{calculators} --- free vs busy slots per pool;
  \item \textbf{activity counters} --- the change counters that timestamp each
        subsystem's last movement.
\end{itemize}

No active scan yields a clear ``No active scan found.''\ message and a non-zero exit.
Repeat it with \cli{watch} for a poor man's dashboard:

\begin{terminal}
\begin{jcmd}
watch -n 5 Jarvis2 monitor
\end{jcmd}
\end{terminal}

\begin{jnote}
The monitor is \textbf{strictly read-only} and cheap by construction. Inside the
runtime, every writer bumps a per-subsystem change counter; the monitoring path
re-reads a subsystem only when its counter moved, and an idle subsystem costs a single
\textsc{get}. You can attach monitors as often as you like without slowing Workers ---
and because the state lives in Redis, the monitor works from any process on the
machine.
\end{jnote}

\section{Reading a snapshot in anger}
\label{sec:monitor-diagnosis}

The snapshot answers the three production questions directly:

\begin{description}[style=nextline, leftmargin=1.2em]
  \item[Is it moving?] \code{completed} grows between two snapshots. If not: are the
        Workers all showing the same current Sample for minutes (a hung calculator ---
        check its \yamlkey{timeout}), or is the task queue empty (sampler finished or
        at an \sampler{AdaptiveBridson} barrier)?
  \item[Where is the bottleneck?] A large task queue with idle calculator slots means
        too few Workers; all slots of one pool busy while Workers wait means that pool
        is the throttle; a growing \emph{archive} queue means storage, not compute, is
        behind (Chapter~\ref{ch:performance}).
  \item[Is anything dying?] The \code{failed} counter and per-Worker status make
        systematic failure visible within seconds of it starting --- long before the
        end-of-run summary.
\end{description}

\section{The \code{[Scan Performance]} block and \code{run\_summary}}
\label{sec:scan-performance}

Every run ends with a two-column console block and the matching
\file{run\_summary.\{json,csv,txt\}} files (Section~\ref{sec:run-summary-file}):

\begin{terminal}
\begin{jout}
[Scan Performance]
  wall_time_sec        :  1284.3
  samples_per_sec      :  0.78
  samples_per_min      :  46.7
  avg_sample_sec       :  1.28
  avg_point_eval_sec   :  9.6
  median_point_eval_sec:  8.9
\end{jout}
\end{terminal}

Two latencies are deliberately distinct:

\begin{itemize}
  \item \code{avg\_sample\_sec} $=$ wall time $/$ finished samples --- the
        \emph{amortized system} latency; parallelism drives it far below any single
        evaluation's cost.
  \item \code{avg\_point\_eval\_sec} / \code{median\_point\_eval\_sec} --- the
        per-point \emph{evaluation} cost as measured inside Workers; this is the
        number that reflects your calculators.
\end{itemize}

The ratio of the two is your effective parallelism --- if it is far from
\yamlkey{workers}, Chapter~\ref{ch:performance} tells you where the loss went.

\begin{jtip}
Keep \file{run\_summary.json} from every production scan (it is tiny). Throughput
regressions after a tool upgrade or a storage migration become one \code{diff} instead
of an argument from memory.
\end{jtip}
